%MB9T5
\setcounter{paragraph}{0}
\setcounter{ex}{0}
\PhanI
\begin{ex}
Dụng cụ nào sau đây \textbf{không} sử dụng trong thí nghiệm giao thoa ánh sáng của Young?
\choice
{Đèn laze}
{Khe cách tử}
{Thước đo độ dài}
{\True Lăng kính}
\loigiai{  }
\end{ex}
\loai{Số vân sáng - tối trong vùng giao thoa}
\begin{ex}
Trong thí nghiệm Young về giao thoa ánh sáng, khoảng cách giữa hai khe là $0,5\ mm$, khoảng cách từ mặt phẳng chứa hai khe đến màn là $2\ m$. Ánh sáng đơn sắc dùng trong thí nghiệm có bước sóng $0,5\ \mu m$. Vùng giao thoa trên màn rộng 26 mm (vân trung tâm ở chính giữa). Số vân sáng là
\choice
{12}
{\True 13}
{14}
{15}
\loigiai{
$i = \dfrac{\lambda D}{a} = \dfrac{0,5 . 10^{-6} . 2}{0,5 . 10^{-3}} = 2 . 10^{-3}$ m $= 2$ (mm)
Số vân sáng :
$N_{s} = 2 \left[ \dfrac{L}{2 i} \right] + 1 = 2 \left[ \dfrac{26}{2 . 2} \right] + 1 = 2 . 6 + 1 = 13$
}
\end{ex}
\begin{ex}
Trong thí nghiệm Young về giao thoa ánh sáng, các khe $S_{1}$, $S_{2}$ được chiếu sáng bởi ánh sáng đơn sắc bước sóng $0,5\ \mu m$. Khoảng cách giữa hai khe $a = 1\ mm$. Khoảng cách giữa mặt phẳng chứa hai khe và màn quan sát E là $D = 3$ m. Xét trong miền giao thoa có bề rộng là $12,75$ mm thì số vân sáng quan sát được là
\choice
{8 vân}
{\True 9 vân}
{12 vân}
{10 vân}
\loigiai{}
\end{ex}

\begin{ex}
Người ta thực hiện giao thoa ánh sáng đơn sắc với hai khe Young cách nhau $0,5\ mm$, khoảng cách giữa hai khe đến màn là $2\ m$, ánh sáng dùng có bước sóng $\lambda = 0,5\ \mu m$. Bề rộng của giao thoa trường là $18\ mm$. Số vân sáng, vân tối quan sát được trên màn là 
\choice
{10 vân sáng, 9 vân tối}
{7 vân sáng, 8 vân tối}
{\True 8 vân sáng, 9 vân tối}
{8 vân sáng, 8 vân tối}
\loigiai{
$i = \dfrac{\lambda D}{a} = \dfrac{0,5 . 10^{-6} . 2}{0,5 . 10^{-3}} = 2 . 10^{-3}$ m $= 2$ (mm)

Số vân sáng :

$N_{s} = 2 \left[ \dfrac{L}{2 i} \right] + 1 = 2 \left[ \dfrac{18}{2 . 2} \right] + 1 = 2 . 4 + 1 = 9$

Số vân tối:

$N_{t} = 2 \left[ \dfrac{L}{2 i} + 0,5 \right] = 2 \left[ \dfrac{18}{2 . 2} + 0,5 \right] = 2 . (4 + 0,5) = 9$
}
\end{ex}
\begin{ex}%[3P5K3-4][Loai3]
Trong thí nghiệm giao thoa ánh sáng với hai khe Young cách nhau 0,5\ mm, ánh sáng có bước sóng 0,5\microm, màn ảnh cách hai khe 2\ m. Bề rộng vùng giao thoa trên màn là 17\ mm. Số vân sáng, vân tối quan sát được trên màn là
\choice
{7 vân sáng, 8 vân tối}
{\True 9 vân sáng, 8 vân tối}
{8 vân sáng, 9 vân tối}
{8 vân sáng, 7 vân tối}
\loigiai{
\begin{itemize}
\item Khoảng vân giao thoa trên màn:
\begin{align*}
i=\dfrac{\lambda.D}{a}=\dfrac{0,5.2}{0,5}=2\ mm\Rightarrow N=\dfrac{L}{2i}=\dfrac{17}{2.2}=4,25\Rightarrow [N]=4.
\end{align*}
\item Suy ra trên màn có \bth{N_s=2[N]+1=2.4+1=9} vân sáng và \bth{N_t=2[N+0,5]=2.4=8} vân tối.
\end{itemize}
}

\end{ex} \haidong{20}
\begin{ex}%[3P5K3-4][Loai3]
Trong thí nghiệm Young về giao thoa ánh sáng, người ta chiếu sáng hai khe hẹp bằng ánh sáng đơn sắc có bước sóng $0,5 \ \mu m$. Khoảng cách giữa hai khe hẹp là 0,5 mm, khoảng cách từ hai khe hẹp đến màn quan sát là 2 m, bề rộng trường giao thoa là 19 mm. Số vân sáng và số vân tối quan sát được trên màn bằng
\choice
{10 vân sáng và 10 vân tối}
{9 vân sáng và 8 vân tối}
{\True 9 vân sáng và 10 vân tối}
{10 vân sáng và 9 vân tối}
\loigiai{\begin{itemize}
\item Khoảng vân: $i=\dfrac{\lambda D}{a}=\dfrac{0,5.2}{0,5}=2\ mm$.
\item Số vân sáng, tối quan sát được trên trường giao thoa:\\
$N_s=1+2.\left[{\dfrac{L}{2i}}\right]=1+2.\left[{4,75}\right]=1+2.4=9$;\
$N_t=2.\left[{\dfrac{L}{2i}+\dfrac{1}{2}}\right]=2.\left[{5,25}\right]=2.5=10$.
\end{itemize}
}

\end{ex} \haidong{20}
\loai{Số vân sáng - tối trên MN cho biết toạ độ}
\begin{ex}%[3P5K3-3][Loai5]
Trong một thí nghiệm giao thoa ánh sáng với khe Young, hai khe cách nhau 2\ mm, khoảng cách từ hai khe tới màn quan sát là 2 m. Ánh sáng đơn sắc có bước sóng $\lambda =0,5\ \mu m$. Cho M và N là hai điểm nằm trong trường giao thoa, chúng nằm khác phía nhau so với vân chính giữa, có $OM=12,3\ mm, ON=5,2\ mm$. Số vân sáng và số vân tối trong đoạn MN là
\choice
{\True 35 vân sáng, 35 vân tối}
{36 vân sáng, 36 vân tối}
{35 vân sáng, 36 vân tối}
{36 vân sáng, 35 vân tối}
\loigiai{
\begin{itemize}
\item Khoảng vân $i=\dfrac{\lambda D}{a}=0,5\ mm$.
\item Vì hai điểm M và N trên màn ở khác phía so với vân sáng trung tâm nên có thể chọn
$x_M=-12,3\ mm$ và $x_N=5,2\ mm$:\\
$\left\{\begin{aligned}& x_M\leq ki=k.0,5\leq x_N\Rightarrow -24,6\leq k\leq 10,4\Rightarrow k=\underbrace{{-24;...;10}}_{{\text{có 35 giá trị}}} \\
& x_M\leq \left({m+0,5}\right)i=\left({m+0,5}\right)0,5\leq x_N\Rightarrow -25,1\leq m\leq 9,9\Rightarrow m=\underbrace{{-25;...;9}}_{{\text{có 35 giá trị}}}
\end{aligned}\right.$
\end{itemize}
}
\end{ex} \haidong{20}
\loai{Số vân sáng - 2 đầu sáng}
\begin{ex}%[3P5K3-4][Loai6]
Trong thí nghiệm giao thoa Young với ánh sáng đơn sắc trên màn chỉ quan sát được 21 vạch sáng và khoảng cách giữa hai vạch sáng đầu và cuối là 40 mm. Tại hai điểm M, N là hai vị trí của hai vân sáng trên màn. Hãy xác định số vân sáng trên đoạn MN biết rằng khoảng cách giữa hai điểm đó là 24 mm.
\choice
{40}
{\True 13}
{41}
{12}
\loigiai{\begin{itemize}
\item Khoảng cách giữa 21 vạch sáng: $\left({21-1}\right).i=40\ mm\Rightarrow i=2\ mm$.
\item Số vân sáng trên MN: $N_s=1+\dfrac{MN}{i}=1+\dfrac{24}{2}=1+12=13$.
\end{itemize}}
\end{ex} \haidong{20}
\begin{ex}%[3P5K3-2][Loai6]
Trong thí nghiệm giao thoa ánh sáng với khe Young, khoảng cách giữa hai khe là 3\ mm, khoảng cách từ hai khe đến màn quan sát là 2 m. Giữa hai điểm P, Q trên màn quan sát đối xứng nhau qua vân sáng trung tâm có 11 vân sáng, tại P và Q là hai vân sáng. Biết khoảng cách PQ là 3\ mm. Bước sóng do nguồn phát ra nhận giá trị
\choice
{$\lambda = 0,65\ \mu m$}
{$\lambda = 0,5\ \mu m$}
{$\lambda = 0,6\ \mu m$}
{\True $\lambda = 0,45\ \mu m$}
\loigiai{
\begin{itemize}
\item Ta có: $i=\dfrac{PQ}{11-1}=0,3.10^{-3} \ m\Rightarrow \lambda =\dfrac{ai}{D}=\dfrac{{3.10}^{-3}{.0,3.10}^{-3}}{2}=0,45.10^{-6}\ m $.
\end{itemize}}
\end{ex} \haidong{20}

\loai{Số vân sáng - 2 đầu tối}
\begin{ex}%[3P5K3-4][Loai7]
Trong thí nghiệm Young về giao thoa ánh sáng, nguồn sáng phát ra ánh sáng đơn sắc, khoảng cách giữa hai khe là 3 mm, khoảng cách từ màn quan sát tới hai khe là 2 m. Giữa hai điểm M, N đối xứng nhau qua vân trung tâm có 13 vân sáng (M và N là 2 vân tối) và $MN = 3,9$ mm. Bước sóng đơn sắc dùng trong thí nghiệm là
\choice
{550 nm}
{520 nm}
{490 nm}
{\True 450 nm}
\loigiai{
\begin{itemize}
\item Ta thấy M, N là 2 vân tối mà giữa MN lại có 13 vân sáng.
\item Vậy MN cũng sẽ có 13 khoảng vân:
$MN=13i=3,9\Rightarrow i=0,3\ mm $ $\Rightarrow \lambda =\dfrac{ia}{D}=\dfrac{0,3. 3}{2}=0,45\ \mu m $.
\end{itemize}
}

\end{ex} \haidong{20}

\loai{Số vân tối - 1 đầu tối - 1 đầu sáng}

\begin{ex}%[3P5K3-4][Loai9]
Trong thí nghiệm Young về giao thoa ánh sáng, trong khoảng rộng 12,5 mm trên màn có 13 vân tối biết một đầu là vân tối còn một đầu là vân sáng. Biết bề rộng trường giao thoa 8,1 mm. Tổng số vân sáng và vân tối có trong miền giao thoa là
\choice
{19}
{\True 17}
{16}
{15}
\loigiai{
\begin{itemize}
\item $12,5i=12,5\ mm\Rightarrow \text{i=1}\ mm\Rightarrow \left\{\begin{aligned}& N_s=2\left[{\dfrac{0,5L}{i}}\right]+1=2\left[{\dfrac{0,5.8,1}{1}}\right]+1=9 \\
& N_t=N_s-1=8
\end{aligned}\right.$
$\Rightarrow N_t+N_s=17$.
\end{itemize}
}

\end{ex} \haidong{20}
\begin{ex}%[3P5K3-4][Loai9]
Trong thí nghiệm Young về giao thoa ánh sáng với ánh sáng đơn sắc, khoảng cách giữa hai khe là 1 mm, khoảng cách từ hai khe tới màn là 2 m. Trong khoảng rộng 12,5 mm trên màn có 13 vân tối biết một đầu là vân tối còn một đầu là vân sáng. Bước sóng của ánh sáng đơn sắc đó là
\choice
{0,48 $\mu m $}
{0,52 $\mu m $}
{\True 0,5 $\mu m $}
{0,46 $\mu m $}
\loigiai{
\begin{itemize}
\item Khoảng cách giữa 2 vân tối liên tiếp là 1i, giữa 1 vân tối và 1 vân sáng liên tiếp là 0,5i.
\item Vậy độ dài khoảng rộng bằng
$12i+0,5i=12,5\Rightarrow i=1\ mm\Rightarrow \lambda =\dfrac{ia}{D}=\dfrac{1. 1}{2}=0,5\ \mu m $.
\end{itemize}
}
\end{ex} \haidong{20}

\loai{Số vân sáng - sáng và ?}
\begin{ex}%[3P5K3-4][Loai10]
Trong thí nghiệm giao thoa ánh sáng với hai khe Young cách nhau 1,8\ mm và cách màn 1,2 m. Ánh sáng đơn sắc làm thí nghiệm có bước sóng 486 nm. Trên bề rộng 3\ mm tính từ vân trung tâm của màn giao thoa quan sát được bao nhiêu vân tối và bao nhiêu vân sáng (không kể vân trung tâm)?
\choice
{10 vân tối, 9 vân sáng}
{9 vân tối, 10 vân sáng}
{\True 9 vân tối, 9 vân sáng}
{8 vân tối, 10 vân sáng}
\loigiai{
\begin{itemize}
\item Khoảng vân giao thoa trên màn: \bth{i=\dfrac{\lambda.D}{a}=\dfrac{0,486.1,2}{1,8}=0,324\ mm}.
\item \bth{N=\dfrac{L}{i}=\dfrac{3}{0,324}=9,3\Rightarrow [N]=9\ \ct{(lẻ)}}.
\item Suy ra trên màn có 10 vân sáng và 9 vân tối.
\end{itemize}
}
\end{ex} \haidong{20}

\begin{ex}
Cho thí nghiệm giao thoa ánh sáng của Young, ánh sáng đơn sắc được dùng trong thí nghiệm có bước sóng là $0,5\  \mu m$. Biết khoảng cách giữa hai khe là $0,5\ mm$, khoảng cách từ hai khe đến màn là $2\ m$. Hai điểm $M, N$ cách vân trung tâm các đoạn $OM = 4,5\ mm$, $ON = 12,5\ mm$ và nằm khác phía so với vân trung tâm. Số vân sáng trên đoạn MN là
\choice
{6}
{7}
{8}
{\True 9}
\loigiai{
$i = \dfrac{\lambda D}{a} = \dfrac{0,5.10^{-6}.2}{0,5.10^{-3}} = 2\ mm$\\
Đoạn MN = $4,5 + 12,5 = 17\ mm \Rightarrow n = \dfrac{17}{2} + 1 = 8,5 + 1 = 9$
}
\end{ex}
\begin{ex}
Trong thí nghiệm Young về giao thoa ánh sáng, giữa hai điểm $P$ và $Q$ trên màn cách nhau $9\ mm$ chỉ có $5$ vân sáng mà tại $P$ là một trong $5$ vân sáng đó, còn tại $Q$ là vị trí của vân tối. Vị trí vân tối thứ $2$ cách vân trung tâm một khoảng là
\choice
{$0,5\ mm$}
{$5\ mm$}
{\True $3\ mm$}
{$0,3\ mm$}
\loigiai{}
\end{ex}
\begin{ex}
Thực hiện giao thoa ánh sáng qua khe Young, biết $a = 0,5\ mm$, $D = 2\ m$. Nguồn S phát ánh sáng trắng gồm vô số bức xạ đơn sắc có bước sóng từ $0,4\ \mu m$ đến $0,76\ \mu m$. Số bức xạ bị tắt tại điểm M trên màn E cách vân trung tâm $0,72\  cm$ và thuộc vùng giao thoa là
\choice
{\True 4}
{5}
{6}
{7}
\loigiai{
M là bức xạ bị tắt nên M là vân tối.
$x_{M} = \left(k + \dfrac{1}{2}\right) i = \left(k + \dfrac{1}{2}\right) \dfrac{\lambda D}{a} \Rightarrow \lambda = \dfrac{2 a . x_{M}}{(2 k + 1) D} = \dfrac{2 . 0,5 . 7,2 . 10^{-3}}{(2 k + 1) 2} = \dfrac{3,6 . 10^{-3}}{2 k + 1}$
$0,4 . 10^{-6} \leq \lambda \leq 0,76 . 10^{-6} \Leftrightarrow 0,4 . 10^{-6} \leq \dfrac{3,6 . 10^{-3}}{2 k + 1} \leq 0,76 . 10^{-6}$
$1,89 \leq 2k + 1 \leq 9 \Rightarrow 0,445 \leq k \leq 4$
$k = 1, 2, 3, 4$
Vậy có 4 bức xạ cho vân tối tại M
}
\end{ex}
\begin{ex}
Trong thí nghiệm Young về giao thoa ánh sáng đơn sắc có bước sóng 600 nm, khoảng cách giữa hai khe là $0,5\ mm$, khoảng cách từ mặt phẳng chứa hai khe đến màn quan sát là $1,5\ m$.
\begin{chc}
Khoảng vân thu được trên màn là 
\choice
{\True 1,8 mm}
{1,4 mm}
{1,2 mm}
{0,6 mm}
\loigiai{  }
\end{chc}
\begin{chc}
Trên màn quan sát, gọi M và N là hai điểm nằm ở hai phía so với vân sáng trung tâm và cách vân sáng trung tâm lần lượt là $6,3\ mm$ và $4,5\ mm$. Số vân sáng nằm trong khoảng MN là
\choice
{\True 6}
{7}
{8}
{9}
\loigiai{  }
\end{chc}
\end{ex}


























\setcounter{ex}{0}
\PhanII
\begin{ex}
\immini[thm]{Trong thí nghiệm Young về giao thoa ánh sáng, hai khe Young $S_1,\ S_2$ được chiếu bằng ánh sáng đơn sắc từ đèn Đ. Khoảng cách giữa hai khe $S_1$,  $S_2$ là $1\ mm$, khoảng cách từ mặt phẳng chứa hai khe đến màn quan sát $E$ là $1,5\ m$. Kết quả, trên màn hứng ảnh thu được $9$ vân sáng, khoảng cách giữa hai vân sáng ngoài cùng là $7,5\ mm$.}{
\includegraphics[width=6 cm]{images/2025_04_24_4fb370520d0729cb661ag-5(1)}}
\choiceTF[t]
{\True Trên màn hứng ảnh, ngay chính giữa miền giao thoa là vân sáng}
{Vân sáng là kết quả của hai sóng ánh sáng gặp nhau triệt tiêu lẫn nhau. Vân tối là kết quả hai sóng ánh sáng ngược pha tăng cường lẫn nhau}
{Khoảng vân là $1,2\ mm$}
{\True Ánh sáng đơn sắc dùng trong thí nghiệm có bước sóng $625\ nm$}
\loigiai{}
\end{ex}

\begin{ex}
Trong thí nghiệm Young về giao thoa ánh sáng, hai khe hẹp được chiếu sáng bởi ánh sáng đơn sắc $\lambda = 0,5\ \mu m$. Khoảng cách giữa hai khe là $a = 1\ mm$. Khoảng cách giữa mặt phẳng hai khe đến màn quan sát là $D = 2,5\ m$.
\choiceTF[t]
{\True Khoảng vân được tính theo công thức: $i = \dfrac{\lambda D}{a}$}
{Khoảng cách giữa 10 vân sáng liên tiếp là $12,5$ mm}
{\True Tại M cách vân sáng trung tâm $5,625$ mm là vân tối thứ 5}
{Nếu bề rộng trường giao thoa (đối xứng) là $13,5$ mm thì số vân sáng quan sát được là 22}
\loigiai{
\begin{itemchoice}
\itemch $i = \dfrac{\lambda D}{a}$.
\itemch Khoảng vân $i = \dfrac{0,5 . 10^{-6} . 2,5}{1 . 10^{-3}} = 1,25 . 10^{-3}$ m $= 1,25$ mm. Khoảng cách giữa 10 vân sáng liên tiếp là $9i = 9 . 1,25 = 11,25$ mm.
\itemch Vân tối thứ 5 ứng với $k=4$, vị trí $x = (4 + 0,5)i = 4,5 . 1,25 = 5,625$ mm.
\itemch Bề rộng $L = 13,5$ mm. Số khoảng vân $N = L/i = 13,5 / 1,25 = 10.8$. Số vân sáng $2N + 1 = 22.6 \approx 23$.
\end{itemchoice}
}
\end{ex}
\begin{ex}
Tiến hành thí nghiệm Young về giao thoa ánh sáng với ánh sáng đơn sắc có bước sóng $0,6\ \mu m$. Khoảng cách giữa hai khe là $0,3\ mm$, khoảng cách từ mặt phẳng chứa hai khe đến màn quan sát là $2\ m$.
\choiceTF[t]
{\True Vân tối thứ 3 cách vân trung tâm $2,5i$}
{\True Hai vân tối liên tiếp cách nhau $4\ mm$}
{Tại điểm $M$ cách vân trung tâm $6\ mm$ là vân sáng bậc 2 tính từ vân trung tâm}
{Hai điểm M và N ở hai phía đối với trung tâm $O$ có $ON=26\ mm$ và $OM=62\ mm$ thì giữa MN có 18 vân tối}
\loigiai{
\begin{itemchoice}
\itemch Khoảng vân $i = \dfrac{\lambda D}{a} = \dfrac{0,6 . 10^{-6} m . 2 m}{0,3 . 10^{-3} m} = 4 . 10^{-3} m = 4 mm$. Vị trí vân tối thứ $k$ là $x_t = (k - 0,5)i$. Vậy vân tối thứ 3 có vị trí $x_{t3} = (3 - 0,5)i = 2,5i$.
\itemch Khoảng cách giữa hai vân tối liên tiếp bằng khoảng vân $i = 4 mm$.
\itemch Vị trí vân sáng bậc $k$ là $x_s = ki$. Với $k=2$, $x_{s2} = 2i = 2 . 4 mm = 8 mm$.
\itemch 
\end{itemchoice}
}
\end{ex}
\begin{ex}
Thí nghiệm Young về giao thoa với ánh sáng, chiếu ánh sáng đỏ có bước sóng $\lambda=0,75\ \mu m$ vào hai khe; khoảng cách giữa hai khe hẹp là $a=1\ mm$, khoảng cách từ hai khe đến màn quan sát là $D=2\ m$.
\choiceTF[t]
{\True Vân sáng trung tâm có màu đỏ}
{\True Khoảng cách giữa hai vân sáng liên tiếp bằng $1,5\ mm$}
{\True Vân tối thứ $5$ cách vân sáng trung tâm đoạn $6,75\ mm$}
{Hai điểm $M$, $N$ ở hai bên vân sáng trung tâm, cách vân sáng trung tâm lần lượt $4,25\ mm$ và $6\ mm$. Giữa $M$ và $N$ có tất cả có $7$ vân sáng}
\loigiai{
\begin{itemchoice}
\itemch
\itemch
\itemch
\itemch
\end{itemchoice}
}
\end{ex}
\setcounter{ex}{0}
\PhanIII

\loai{Trả lời ngắn}

\begin{ex}
Trong thí nghiệm khe Young về ánh sáng đơn sắc, khoảng cách giữa 7 vân sáng liên tiếp trên màn là $21,6 ~mm$. Trong vùng giao thoa trên màn rộng $31 ~mm$ vân trung tâm ở chính giữa. Tìm số vân sáng quan sát được.
\shortans[oly]{9}
\loigiai{
\begin{itemize}
\item Khoảng cách 7 vân sáng liên tiếp: $6 i=21,6 ~mm \Rightarrow i=3,6 ~mm$
\item Số vân sáng: $-\dfrac{L}{2 i} \leq k \leq \dfrac{L}{2 i} \Rightarrow-\dfrac{31}{2 \cdot 3,6} \leq k \leq \dfrac{31}{2 \cdot 3,6}$
\end{itemize}
Đáp số: 9
}
\end{ex} \haidong{20}
\begin{ex}
Thực hiện giao thoa ánh sáng bằng khe Young với ánh sáng đơn sắc có bước sóng là $\lambda$. Người ta đo khoảng cách giữa vân sáng và vân tối nằm cạnh nhau là $1 ~mm$. Trong khoảng giữa hai điểm M và N trên màn và ở hai bên so với vân trung tâm, cách vân này lần lượt là $6 ~mm$ và $7 ~mm$ có bao nhiêu vân sáng?\shortans[oly]{6}
\loigiai{
\begin{itemize}
\item Khoảng cách giữa vân sáng và vân tối cạnh nhau: $0,5 i=1 \Rightarrow i=2 ~mm$
\item Số vân sáng: $\dfrac{x_{M}}{i}<k<\dfrac{x_{N}}{i} \Rightarrow-\dfrac{6}{2}<k<\dfrac{7}{2} \Rightarrow-3<k<3,5$
\end{itemize}
Đáp số: 6
}
\end{ex} \haidong{20}
\begin{ex}
Trong thí nghiệm về giao thoa ánh sáng của Young, khoảng cách giữa hai khe là $1 \  mm$, khoảng cách từ hai khe đến màn là $1 \   m$, bước sóng ánh sáng là $0,5 \  \mu m$. Xét hai điểm M, N ở cùng một phía đối với vân sáng trung tâm có tọa độ lần lượt là $1 \  mm$ và $5,25 \  mm$. Trên đoạn MN có bao nhiêu vân sáng?
\shortans[oly]{9}
\loigiai{

}
\end{ex}
\begin{ex}
Trong thí nghiệm Young về giao thoa ánh sáng hai khe cách nhau $0,4\ mm$, khoảng cách từ hai khe đến màn là $1\ m$. Ánh sáng đơn sắc dùng trong thí nghiệm có bước sóng $0,5\ \mu m$. Vùng giao thoa trên màn có bề rộng $13\ mm$ đối xứng qua vân sáng trung tâm chứa bao nhiêu vân sáng?
\shortans[oly]{11}
\loigiai{

}
\end{ex}
\begin{ex}
Trong thí nghiệm Young về giao thoa ánh sáng, hai khe được chiếu bằng ánh sáng đơn sắc có bước sóng $\lambda=640\ nm$. Khoảng cách giữa hai khe là $a=1\ mm$, khoảng cách từ mặt phẳng chứa hai khe đến màn quan sát là $D=100\ cm$. Bề rộng miền giao thoa là $L=1,5\ cm$. Tổng số vân tối có trong miền giao thoa là bao nhiêu?
\shortans[oly]{24}
\loigiai{

}
\end{ex}
\begin{ex}
Thí nghiệm Young về giao thoa với ánh sáng, chiếu ánh sáng đỏ có bước sóng $\lambda=0,75\ \mu m$ vào hai khe; khoảng cách giữa hai khe hẹp là $a=1\ mm$, khoảng cách từ hai khe đến màn quan sát là $D=2\ m$. Hai điểm $M$, $N$ ở hai bên vân sáng trung tâm, cách vân sáng trung tâm lần lượt $4,25\ mm$ và $6\ mm$. Giữa $M$ và $N$ có tất cả có bao nhiêu vân sáng? 
\shortans[oly]{7}
\loigiai{
}
\end{ex}
















